\section{Introduction}
\label{sec:introduction}

Coastal lagoons are among the most productive coastal environments on Earth,
concentrating fisheries and biodiversity along roughly 13\% of the global
shoreline \citep{Carrasco2016, Kjerfve1994}. The shallow, semi-enclosed
geometry that sustains this productivity also makes them acutely sensitive to
the atmospheric and tidal forcing that drives water exchange and renewal. In
the Mediterranean, where tidal amplitudes rarely exceed 0.3~m, the balance
among tide, wind stress and seiches varies from one lagoon to another, but wind
consistently emerges as a significant driver of horizontal exchange, and
through it of the residence times and salinity gradients that condition benthic
habitat \citep{Niedda2007, Umgiesser2014}. In many of these lagoons that
habitat is seagrass, principally \emph{Posidonia oceanica} (L.) Delile and
\emph{Cymodocea nodosa} (Ucria) Ascherson, and the relation runs in both
directions, since the canopies that depend on the circulation also act on it,
attenuating incident wave energy \citep{Infantes2012} and extracting momentum
from the mean flow through a drag that modifies near-bed turbulence and the
vertical velocity structure \citep{Nepf2012}. These meadows are declining,
global seagrass extent having fallen in all seven bioregions, by 19.1\% of the
surveyed area since 1880 \citep{Waycott2009, Dunic2021}, with the fastest
recent losses where meadows are exposed to coastal development and degraded
water quality \citep{Turschwell2021}. How circulation sets residence times,
dispersal and habitat connectivity in vegetated lagoons is therefore a question
for management as much as for science, and it concerns transport pathways, a
quantity that water level, the usual validation target of lagoon models, does
not measure directly.

The \emph{Stagnone di Marsala}, on the north-western coast of Sicily, is a
lagoon of this kind. Its mean depth is about 0.95~m, and its floor is largely
covered by meadows of \emph{P.\ oceanica} and \emph{C.\ nodosa}
(Section~\ref{sec:study_area}). Satellite monitoring over 2003--2024 has
documented a 75\% reduction in the extent of the \emph{P.\ oceanica} atoll and
cordon formations \citep{Maltese2025}, so the vegetated state a model of this
system must represent is itself changing.

Numerical studies of the Stagnone have each been built around the question they
addressed. \citet{LaLoggia2004} and \citet{DiMarca2009} used depth-averaged
models to characterise tidal flushing, residence times and turbidity, for which
the basin-scale water budget is the relevant quantity, and
\citet{Ciraolo2009a} showed through particle tracking how strongly simulated
trajectories depend on the velocity field they are advected through.
\citet{DeMarchis2012} resolved the vertical and showed that wind drives a
recirculation with a near-bed return flow even at sub-metre depths, with the
meadow represented through Nikuradse roughness heights assigned by seagrass
class and the model validated against velocity and water level at several
stations. \citet{Ingrassia2024} calibrated a depth-averaged model against water
level for a range of vegetation scenarios, a setting in which a uniform
equivalent roughness is the appropriate representation of the canopy. Building
on these studies, three aspects of the problem remain open for this basin. The
first is the representation of the canopy as a drag resolved over the vertical
extent of the plants, a formulation that acts on the shape of the velocity
profile as well as on its magnitude \citep{Nepf2012} and that is now available
in three-dimensional solvers. The second is the effect of waves, which in
comparably shallow Mediterranean lagoons alter bed stress and mobilise sediment
\citep{Carniello2011, Morote2025}, together with that of a bed evolving under
their combined action with the currents. The third is the assessment of
transport itself, the models having been evaluated mainly on water level and on
currents at fixed stations. Whether a model that reproduces water level also
reproduces transport, and which of the processes it carries are responsible
when it does, therefore remains to be established for this basin, and the
question extends to shallow vegetated lagoons in general, where Lagrangian
observations are known to constrain a model differently from water level
\citep{Liu2011, Revelard2021}.

The study takes up these three aspects by separating, in a three-dimensional model of
the Stagnone, the contributions of the canopy, of wave coupling and of bed
mobility, and by scoring every combination of the three against both water
level and observed drifter trajectories (Section~\ref{sec:data_methods}). It
pursues three objectives. The first is to rank the three processes by their
effect on Lagrangian transport skill. The second is to establish whether water
level, the quantity against which shallow lagoon models are conventionally
validated, can identify that ranking. The third is to measure how far the
apparent effect of each process depends on how the other two are represented.
The second and third are not specific to this basin, concerning what a class of
validation data can certify and how a factorial behaves when a dominant factor
is left out of it.
